\chapter{Marco Metodológico}
Para el desarrollo de este proyecto se seguirá una metodología basada en el enfoque ágil basada en el proceso unificado de software (Agile Unified Process), que permite una adaptación continua a los cambios y una entrega incremental de funcionalidades.\\
Este enfoque se caracteriza por las siguientes fases:
\begin{itemize}
    \item \textbf{Inicio:} En esta fase se definirán los objetivos del proyecto, se identificarán los requisitos iniciales y se establecerá el alcance del sistema.
    \item \textbf{Elaboración:} Durante esta fase se realizará un análisis detallado de los requisitos, se diseñará la arquitectura del sistema y se planificarán las iteraciones de desarrollo.
    \item \textbf{Construcción:} En esta fase se llevará a cabo el desarrollo del sistema en iteraciones, implementando las funcionalidades definidas en cada ciclo.
    \item \textbf{Transición:} Finalmente, en esta fase se realizará la entrega del sistema, se llevarán a cabo pruebas finales y se proporcionará soporte para la implementación.
\end{itemize}

\section{Recolección de Datos}
Para la recolección de datos se utilizarán las siguientes técnicas:
\begin{itemize}
    \item \textbf{Encuestas:} Se diseñarán encuestas para recopilar información sobre las necesidades y expectativas de los usuarios en relación con la gestión de notas académicas.
    \item \textbf{Entrevistas:} Se llevarán a cabo entrevistas con estudiantes y profesores para obtener una comprensión más profunda de sus experiencias y desafíos actuales.
    \item \textbf{Análisis de Documentos:} Se revisarán documentos y materiales académicos para identificar patrones y requisitos específicos relacionados con la gestión de notas.
\end{itemize}   

\section{Herramientas}

Para el desarrollo del proyecto se utilizarán las siguientes herramientas:
\begin{itemize}
    \item \textbf{Lenguajes de Programación:} Python para el desarrollo del backend y JavaScript para el frontend.
    \item \textbf{Frameworks:} Django para el backend y React para el frontend.
    \item \textbf{Bases de Datos:} PostgreSQL para el almacenamiento de datos.
    \item \textbf{Control de Versiones:} Git y GitHub para la gestión del código fuente.
    \item \textbf{Herramientas de Diseño:} Figma para la creación de prototipos de interfaz de usuario.
\end{itemize}

\section{Scrum}
\noindent Para la planeacion y desarrollo de la aplicación Gestor inteligente para notas academicás se usara el metodo agil \textit{Scrum}. El cual podemos entender como una metodología ágil de desarrollo que ayuda a los equipos a desarrollar productos en ciclos cortos llamados \textit{sprints}, permitiendo obtener retroalimentación rápida, adaptarse a cambios y mejorar continuamente. \textit{Scrum} fomenta la autoorganización del equipo, la entrega de valor frecuente y la colaboración directa para abordar problemas complejos de forma iterativa. Su nombre proviene del rugby, simbolizando el trabajo conjunto del equipo para alcanzar un objetivo común. \cite{Scrum}
